\begin{table}[htbp]
\centering
\caption{Fresh actual-cost comparison with a nonuniform conventional method}\label{tab:adaptive54}
{\small\setlength{\tabcolsep}{4pt}
\begin{tabular}{rrcc}
\hline
Method & Pass & Actual expected cost & Gain from initial policy \\
\hline
W & 0 & [0.60337, 0.62425] & -- \\
W & 1 & [0.59714, 0.61789] & [0.00129, 0.01130] \\
W & 2 & [0.59173, 0.61245] & [0.00662, 0.01681] \\
F & 0 & [0.59081, 0.61156] & -- \\
F & 1 & [0.59065, 0.61141] & [-0.00462, 0.00494] \\
F & 2 & [0.59065, 0.61141] & [-0.00462, 0.00494] \\
A & 0 & [0.59068, 0.61143] & -- \\
A & 1 & [0.59060, 0.61135] & [-0.00470, 0.00486] \\
A & 2 & [0.59060, 0.61135] & [-0.00470, 0.00486] \\
\hline
\end{tabular}}
\par\smallskip
{\footnotesize A is surplus-driven FVI; W and F denote the same generators as above. All three are reconstructed in the same separately frozen cohort and receive two full improvement passes. W and F policy identities reproduce the primary cohort, while the cost paths are a fresh family. Nonuniformity and pilot work are recorded explicitly.}
\end{table}
